\vpage{1}{Three objects, three different questions}{PAPER V / A WORKED VISUAL COMPANION}{Chapter V: Section 1, Section 6, Section 11, Section 13.}
A packet calculus can reorganize a signal exactly, choose among finitely many representations, or describe a nonlinear coordinate geometry. These are not the same task. This guide keeps one concrete object in view for each, so that the reader can compute rather than only follow arrows.
\vfig{three_worlds}
\head{A signal that can be checked by hand}
Our main finite signal is $f=s+j\in\QQ^{16}$. The background repeats
\[
\tfrac14(0,1,2,1,0,-1,-2,-1),
\]
and $j[n]=3$ at $6\le n<12$, with $j[n]=0$ elsewhere. This is a discrete background model, not an assertion of continuum smoothness. Its energy is exactly $57$. We use it for residue coordinates, Haar transforms, coefficient deletion, and known-packet cancellation.
\head{A tree whose alternatives can all be seen}
The full binary tree of depth three has only $26$ complete prunings. We enumerate all of them, compute a unique mixed-depth optimum, and certify its stability with rational costs. The displayed objective is deliberately explicit, not claimed to be a universal signal-quality score.
\head{A coordinate orbit with no hidden numerical poles}
The four reflection points $\Gamma=i^k/2$ map to four exact impedances. A second orbit passes through infinity; homogeneous coordinates make that point an ordinary part of the construction.
\takeaway{Exact coordinates preserve information. An optimizer chooses coordinates. An estimator must still identify information not already supplied.}

\vpage{2}{An atlas of the objects}{READING CONVENTIONS / KEEP THE TYPES VISIBLE}{Chapter V: Section 20 (notation crosswalk), Section 21 (formal correspondence).}
Several symbols in the paper describe different realizations of the same operator language. The following distinctions are used throughout this guide.

\begin{tabularx}{\linewidth}{@{}p{.17\linewidth}YY@{}}\toprule
Object & Literal meaning & What it is not\\\midrule
$S,\delta_0$ & Sample shift and unit impulse & A finite truncation of every $\ell^2$ signal\\
$P_{N,r}f$ & A masked signal on the original lattice & The shorter decimated channel\\
$f_r[q]$ & The value $f[Nq+r]$ & An average of neighboring samples\\
$D$ & One channel-index delay & Necessarily one original-sample delay\\
$E(D),R(D)$ & Finite Laurent analysis and synthesis matrices & Values at one test frequency\\
$W_\alpha$ & The addressed Haar analysis operator & An edge detector by itself\\
$W_\alpha^*W_\alpha$ & A projector on the original signal space & A scalar packet energy\\
$\Lambda$ & A complete covering antichain of a finite tree & An arbitrary collection of nodes\\
$j,\widehat j$ & Known and estimated discontinuity packets & Interchangeable hypotheses\\
$\Gamma,z$ & Reflection coefficient and normalized impedance & Coefficients that can be linearly interchanged\\\bottomrule
\end{tabularx}
\head{Three arithmetic conventions}
FIR examples use bilateral sequences with finite support and \emph{zero extension}. Haar examples use vectors of length $16$, with lengths divisible by the required powers of two. The three-tap smoother and nonlinear neighbor example use \emph{periodic} indexing. These boundary conventions are stated where they enter. None is silently substituted for another.
The polynomial variable $t$ on \pagerefV{20} indexes coefficients. The delay symbol $D$ acts on channels. The Smith variable $z$ is impedance. Keeping them separate prevents a visual resemblance from becoming an invalid identity.
\takeaway{A correct formula needs the correct object, product, norm, and boundary convention. A familiar-looking matrix does not supply them automatically.}

\vpage{3}{Decimation is a permutation, not a loss}{EXACT ADDRESSES / THE RUNNING SIGNAL}{Chapter V: Theorem 2.1, Corollary 2.2.}
For four phases, write every sample address uniquely as $n=4q+r$ with $0\le r<4$. Analysis stores $f[n]$ at channel address $(r,q)$. Synthesis puts that exact value back at $4q+r$.
\vfig{addresses}
\platecaption{Every cell contains an original address and its rational sample value. Shading marks the known jump packet, not an estimated feature or a frequency band.}
\head{What has changed? Only where the values are stored}
The complete tuple of channels contains sixteen values, just as the original vector does. A single decimated channel is shorter, but keeping all four loses no samples. The residue projector $P_{4,r}f$ is different: it keeps the original lattice and writes zeros outside phase $r$.
\[
\mathcal S_4\mathcal A_4f=f,
\qquad
\|f\|_2^2=\sum_{r=0}^3\|f_r\|_2^2=57.
\]
The norm on the channel tuple is the sum-of-squares Hilbert norm, not the maximum of the channel norms. This distinction matters in the infinite-dimensional statement.
\head{Negative addresses are ordinary addresses}
Euclidean division, rather than truncation toward zero, gives
$-1=4(-1)+3$ and $-5=4(-2)+3$. The finite tests exercise both positive and negative sample addresses. The algebraic identity applies before any plotting window or finite buffer is chosen.
\takeaway{Downsampling one channel discards information. Splitting into all residue channels is an invertible address change.}

\vpage{4}{Refine packets on a common lattice}{EXACT REFINEMENT / WHAT CAN BE ADDED}{Chapter V: Theorem 2.3.}
An even sample belongs to exactly one of residues $0$ and $2$ modulo four. Therefore the coarse even mask is the sum of two disjoint fine masks:
\[
P_{2,0}=P_{4,0}+P_{4,2}.
\]
\vfig{refinement}
\platecaption{The common horizontal axis is the original sample index. Empty cells denote masked-out addresses. There are no interpolated samples between the displayed cells.}
\head{Why the coordinate convention matters}
The decimated channel $f_0[q]=f[2q]$ has its own index $q$. At modulus four the descendants are $f[4q]$ and $f[4q+2]$. These shorter arrays must be interleaved to recover the coarse channel. Adding them entrywise instead produces sums of different original samples.
For example, at the first descendant address the values $f[0]=0$ and $f[2]=1/2$ are two separate values. Their sum $1/2$ cannot tell us both of them. Exact refinement is a disjoint partition, not an averaging rule or a new measurement.
\head{The general formula}
When $N$ divides $M$,
\[
P_{N,r}=\sum_{\substack{0\le a<M\\a\equiv r\pmod N}}P_{M,a}.
\]
The divisibility hypothesis ensures that every fine residue class lies entirely inside one coarse class. Without it, this particular descendant relation is not the right refinement rule.
\takeaway{Finer coordinates can preserve the object exactly. Compressing their values to averages is an additional operation, with an additional information-loss question.}

\vpage{5}{A filter becomes a matrix of actual delays}{FIR TRANSFER / A FOUR-TAP EXAMPLE}{Chapter V: Theorem 3.1.}
Take the causal filter $h=(1,2,-1,1)$, with taps at $0,1,2,3$, acting by $y[n]=\sum_qh[q]f[n-q]$. Resolve the input into even and odd channels. Let $D$ denote one delay in the \emph{channel} index.
\vfig{fir_transfer}
\head{Reading one matrix entry}
The odd-input contribution to the even output is $(2D+D^2)f_1$. The tap at original delay one reaches the preceding odd channel sample; the tap at delay three reaches the odd channel two positions back. These are different delays despite coming from the same input phase.
The two channel equations are
\[
\begin{aligned}
y_0[n]&=f_0[n]-f_0[n-1]+2f_1[n-1]+f_1[n-2],\\
y_1[n]&=2f_0[n]+f_0[n-1]+f_1[n]-f_1[n-1].
\end{aligned}
\]
Interleaving gives exactly the direct convolution, including its boundary samples under zero extension.
\head{The general address calculation}
For $N$ phases, the block from input $s$ to output $k$ has kernel
$h_{ks}[m]=h[Nm+k-s]$. This comes from the exact identity
$Nn+k-q=N(n-m)+s$. It is a finite reindexing of taps, so no numerical transform of the input is needed. Evaluation of a bilateral transform is a separate step with the summability hypotheses stated in the manuscript.
\takeaway{The matrix does not approximate the filter. Its entries record which input phase is used, how far back to read, and what coefficient to multiply.}

\vpage{6}{One sample is not one channel step}{DELAY UNITS / AN IMPULSE MAKES THE DIFFERENCE VISIBLE}{Chapter V: Theorem 4.1.}
A one-sample delay swaps the two phases and wraps one of them into the preceding channel index. Its matrix is $J(D)$, not $DI_2$.
\vfig{sample_delay}
\head{Trace one impulse}
The input $\delta_0$ has a single nonzero value in even channel address zero. Under $J(D)$ it appears in odd channel address zero: this is the original sample $\delta_1$. Under $DI_2$ it stays even but moves to channel address one: this is $\delta_2$.
Thus
\[
J(D)^2=DI_2,
\qquad
H(D)=I_2+2J(D)-J(D)^2+J(D)^3
\]
recovers the previous page's four-tap transfer matrix. The same delay algebra explains both the off-diagonal phase swap and the channel powers.
\head{Signed delays need the same care}
For an original-sample delay $d$ at modulus $N$,
\[
y_k[n]=f_{(k-d)\bmod N}
\!\left[n+\left\lfloor\frac{k-d}{N}\right\rfloor\right].
\]
This includes advances when $d<0$. The finite verification checks the formula on signed sample addresses. A bank with product $D^qI_N$ delays the original signal by $Nq$ samples, not $q$ samples.
\takeaway{Always state the unit of a delay. A phase permutation and a channel delay together can represent one original sample.}

\vpage{7}{Perfect reconstruction is an operator identity}{FILTER BANKS / WHY A FREQUENCY SPOT-CHECK IS INSUFFICIENT}{Chapter V: Theorem 4.1.}
Consider the finite Laurent bank
\[
E(D)=\mat{1&D\\0&1},\qquad
R(D)=\mat{1&-D\\0&1}.
\]
Its synthesis cancels the exact delayed cross-channel contribution introduced by analysis.
\vfig{perfect_reconstruction}
\head{Two inverses, all inputs}
Direct polynomial multiplication gives $RE=ER=I_2$. Since $D$ acts as a shift, these identities imply perfect reconstruction on every bilateral channel sequence. With finite kernels, composition requires only finitely many operations at each address.
The converse uses impulses. Insert an impulse in one input channel. Each output channel lists every coefficient of the corresponding transfer polynomial. If all such responses vanish for $RE-I$, every coefficient is zero; hence the polynomial matrix itself is zero.
\head{A deliberately misleading test}
The matrix $B(D)=\operatorname{diag}(D,1)$ satisfies $B(1)=I_2$. At that one transform value it appears perfect. In the time domain, however, it delays only the first channel. An impulse in that channel immediately distinguishes it from the identity.
This is why a small set of successful frequency evaluations is not automatically a proof of the full Laurent identity. A separate degree bound and a sufficient interpolation argument would be needed to turn finitely many evaluations into such a certificate.
\takeaway{Exact reconstruction concerns every coefficient of the transfer operator, not just an attractive response curve or one evaluation point.}

\vpage{8}{A jump can hide from the first Haar split}{HAAR ENERGY / EDGE ALIGNMENT}{Chapter V: Corollary 4.2.}
For the running signal, Haar analysis gives
\[
a[n]=\frac{f[2n]+f[2n+1]}{\sqrt2},\qquad
d[n]=\frac{f[2n]-f[2n+1]}{\sqrt2}.
\]
The plot scales both channels by $\sqrt2$ so that all displayed ordinates are rational.
\vfig{haar_alignment}
\head{The exact energy calculation}
The low channel has energy $227/4$, and the detail channel has energy $1/4$. Their sum is $57$, exactly the input energy. The normalization is essential: unnormalized sums and differences would multiply the total energy by two.
\head{Why is the first detail channel so small?}
The jump starts at sample $6$ and ends just before sample $12$. Both lie \emph{between} the pairs $(0,1),(2,3),\ldots,(14,15)$. Inside every pair the jump contributes either $(0,0)$ or $(3,3)$, so its first-level difference is zero. The small detail values come entirely from the discrete background.
This does not mean the edge is absent from the transform. The next split of the low channel compares neighboring pairs, and those comparisons do cross the jump. The next pages expose that contribution explicitly.
\takeaway{Haar energy is exact, but a particular detail level is not a universal edge detector. Edge visibility depends on alignment, scale, and the chosen packets.}

\vpage{9}{Lifting preserves information, not energy}{PREDICT--UPDATE / INVERTIBILITY VERSUS ISOMETRY}{Chapter V: Theorem 4.3.}
For an even/odd pair $(e,o)$, form the prediction residual $d=o-e$, then update $s=e+d/2$. The inverse reads $e=s-d/2$, followed by $o=d+e$.
\vfig{lifting_geometry}
\platecaption{The square $[0,4]^2$ and its exact image under the lifting matrix have the same area. The separate pair $(2,6)$ illustrates that Euclidean energy need not be preserved.}
\[
\mat{s\\d}=\mat{1/2&1/2\\-1&1}\mat{e\\o},
\qquad
\mat{e\\o}=\mat{1&-1/2\\1&1/2}\mat{s\\d}.
\]
\head{What determinant one does establish}
The matrix is invertible, and in these real coordinates it preserves oriented area. The predict and update operations can be undone in reverse order. Applied to each pair, the construction gives exact finite-signal reconstruction.
\head{What determinant one does not establish}
The pair $(2,6)$ becomes $(4,4)$. Its squared Euclidean norm changes from $40$ to $32$. The lifting matrix is not unitary. Energy accounting must therefore use its actual coordinate metric or a separately normalized transform; invertibility alone cannot replace the Haar energy identity.
All values in this example are rational. Integer-to-integer lifting with a rounding rule is a different algorithm. Its inverse and error conventions require their own specification and are not silently inferred from the displayed rational matrix.
\takeaway{Perfect reconstruction, determinant one, and norm preservation are different properties. This transform has the first two, not the third.}

\vpage{10}{The second level resolves the hidden change}{ACTUAL WAVELET PACKETS / FOUR RATIONAL CHANNELS}{Chapter V: Theorem 5.1.}
Apply Haar analysis to \emph{both} children. A binary address records the chronological low/detail choices, with $0$ for low and $1$ for detail. Each depth-two channel contains four coefficients.
\vfig{depth_two}
\head{The edge that first-level detail missed}
Packet $01$ is the detail child of the first low channel. It contains $-11/4$, the large change between neighboring pair averages. Packet $10$ vanishes for this particular signal because the first-level differences alternate in canceling pairs. Neither feature is a universal identity for all inputs.
The packet energies are
\[
49,\qquad 31/4,\qquad 0,\qquad 1/4,
\]
which again sum to $57$. Two factors of $1/\sqrt2$ explain why the depth-two coefficients are rational halves and quarters.
\head{Coefficients and projectors are different objects}
$W_\alpha f$ is the short coefficient vector. The original-space component is $W_\alpha^*W_\alpha f$. Summing these synthesized components, not the short arrays entrywise, reconstructs $f$:
\[
f=\sum_{|\alpha|=2}W_\alpha^*W_\alpha f.
\]
The finite verification constructs the actual analysis matrices and checks their product with their adjoints. The paper's $\ell^2$ statement additionally establishes boundedness and the existence of the literal Hilbert adjoints.
\takeaway{A packet address specifies an operator word. To add different packets as signals, synthesize them back into the common signal space first.}

\vpage{11}{Deleting coefficients is the approximation step}{SPARSITY / AN EXACT ERROR LEDGER}{Chapter V: Section 10.}
The complete depth-two transform is exact. Approximation begins only when we discard some of its sixteen coefficients. Order them by squared magnitude, with a deterministic address-and-position rule for ties.
\vfig{sparse_tail}
\head{Every discarded square has a meaning}
In this fixed orthonormal basis, the squared reconstruction error equals the sum of the squared discarded coefficients. For our running signal:
\begin{center}
\begin{tabular}{r|rrrrrr}
kept coefficients $k$ & 1&2&4&8&12&16\\\hline
squared error & $59/4$&$115/16$&$11/16$&$1/4$&$0$&$0$
\end{tabular}
\end{center}
The four zero coefficients belong to packet $10$, so twelve coefficients suffice for exact reconstruction here. This is a property of the selected signal and basis, not a guarantee that all edge signals need only twelve coefficients.
\head{Do not compare incompatible budgets}
A cutoff $m$ in a one-dimensional signed Fourier spectrum means $2m+1$ bins, including zero. A budget of $k$ largest coefficients is a different rule. Real Fourier synthesis may also require conjugate-pair closure. Neither count alone equals a bit-storage budget once amplitudes, phases, indices, or a selected tree must be encoded.
The manuscript's Fourier-versus-Haar experiment is a finite comparison, not a proof that one basis dominates every signal class. Here we isolate the part that can be checked exactly: reconstruction error within one fixed finite basis.
\takeaway{Changing coordinates creates no error. Throwing away coefficients does, and orthogonality makes that error exactly accountable.}

\vpage{12}{A mixed-depth basis must cover every branch}{COMPLETE PRUNINGS / NO GAPS AND NO DOUBLE COUNTING}{Chapter V: Definition 11.1, Theorem 11.2.}
A complete pruning chooses nodes of a finite refinement tree so that every root-to-terminal path meets exactly one selected node. The horizontal positions below are \emph{terminal packet addresses}, not time intervals in the original signal.
\vfig{complete_cuts}
\head{Three candidate selections}
The complete selection $\{0,10,11\}$ covers every terminal branch once. The selection $\{0,00,1\}$ covers the $00$ branch twice: it is not an antichain. The selection $\{00,01,10\}$ misses the $11$ branch entirely.
For actual Haar packets, set $P_v=W_v^*W_v$. Parent-to-child conservation gives
\[
P_v=P_{v0}+P_{v1}.
\]
Replacing unselected parents by their children until the complete pruning is reached proves $\sum_{v\in\Lambda}P_v=I$. This is why the coefficients at different depths still form an exact coordinate system.
\head{An exhaustive finite check}
At depth three the pruning count obeys $a_0=1$ and $a_{d+1}=1+a_d^2$, giving $1,2,5,26$. The supplied suite checks all $26$ covering conditions and all $26$ actual $16\times16$ projector-sum identities. These finite checks illustrate, but do not replace, the universal tree theorem.
\takeaway{Mixed depth is allowed. Missing a terminal path or selecting an ancestor and its descendant is not.}

\vpage{13}{Solve every alternative without searching every alternative}{ADDITIVE COSTS / EXACT DYNAMIC PROGRAMMING}{Chapter V: Theorem 11.3, Theorem 11.4.}
Give every node the rational stop cost shown below. A node may either stop and pay its own cost, or refine and pay the already optimized costs of its children. The shaded nodes form the selected complete pruning.
\vfig[.94]{dynamic_program}
\[
V(v)=\min\{c_v,\ V(v0)+V(v1)\},
\qquad V(v)=c_v\text{ at a terminal node}.
\]
\head{Read the calculation from the bottom upward}
At $10$, the children cost $1+1=2$, better than stopping for $7$. At $11$, stopping costs $4$, better than $2+3=5$. The right subtree therefore costs $6$, better than its stop cost $8$. The left subtree costs $1+2=3$, better than $5$. At the root, $3+6=9$ beats $20$.
The unique optimum is
\[
\Lambda_* = \{00,01,100,101,11\},\qquad J(\Lambda_*)=9.
\]
All $26$ complete prunings are separately enumerated and give the same minimum. At equality, the rule stops, making the selector deterministic. A singleton tree needs no comparison.
The costs here are supplied numbers. The theorem optimizes that stated additive objective; it does not prove those numbers encode the best denoising criterion. Coupled costs that cannot be written as a sum of selected-node costs need a different optimizer.
\takeaway{Global optimality follows from the exact stop-or-refine decomposition of a complete pruning, not from a greedy assumption.}

\vpage{14}{How far can the costs move before the tree changes?}{A GLOBAL MARGIN / NOT A LOCAL HEURISTIC}{Chapter V: Theorem 11.3, Corollary 11.5.}
The best pruning on the previous page costs $9$; the cheapest competitor costs $10$. Its exact global gap is therefore $\Gamma=1$. The plate lists every complete-pruning cost in sorted order, giving each pruning its own horizontal rank.
\vfig{selector_margin}
\head{Propagate an error budget as carefully as a cost}
Suppose every node cost changes by at most $\varepsilon_v=1/32$. A pruning may select as many as eight leaves, so its total cost can change by at most $8/32=1/4$. The recursive envelope proves this without guessing which pruning will be selected:
\[
E(v)=\max\{\varepsilon_v,\ E(v0)+E(v1)\},\qquad E(\varnothing)=\tfrac14.
\]
For any competitor $\Lambda$, the changed cost difference is bounded below by
\[
J'(\Lambda)-J'(\Lambda_*)\ge \Gamma-2E(\varnothing)=\tfrac12>0.
\]
Thus the optimum is unchanged for \emph{every} permitted simultaneous perturbation, not just a sampled noise realization. The envelope is over all feasible prunings, not just the five nodes in the current optimum.
\head{Strict inequality has content}
If two competing costs are $9$ and $10$, moving them to $9.5$ and $9.5$ removes uniqueness. Equality in $2E=\Gamma$ is therefore not a unique-selector certificate. For ordinary probability, the random node costs and the joint good event must be measurable. Separate nodewise probability bounds do not automatically give the probability of that joint event.
\takeaway{The exact selector is stable because its global gap exceeds the total worst-case movement of both competing costs.}

\vpage{15}{Retain the jump and approximate the other packet}{AN EXACT IDENTITY / A CONDITIONAL INPUT}{Chapter V: Theorem 6.1, Corollary 6.2.}
Return to the running signal $f=s+j$. For a concrete finite calculation, use the periodic three-tap smoother
\[
(Qx)[n]=\tfrac14x[n-1]+\tfrac12x[n]+\tfrac14x[n+1],\qquad n\in\mathbb Z/16\mathbb Z.
\]
The plots compare smoothing the full signal with retaining its supplied jump packet. Every plotted sample is a rational number computed from the displayed formula.
\vfig{known_edge}
\head{The cancellation is the theorem}
Define $\widetilde f=Qs+j$. Then, sample by sample,
\[
\boxed{f-\widetilde f=s-Qs.}
\]
By contrast, linearity gives
\[
f-Qf=(s-Qs)+(j-Qj).
\]
The second error includes a new term caused by smoothing the discontinuity. In this finite example that term is concentrated near the two jump boundaries. The identity itself requires no linearity of the approximation applied to $s$; linearity is used only in the comparison with $Qf$.
\head{The input is not discovered by subtraction}
The data here include the correct $j$. Replacing it with a wrong $\widehat j$ retains the error $j-\widehat j$, even when $s$ is approximated perfectly. Multiple amplitudes and overlapping shapes are allowed if their full sum is supplied, but the identity does not estimate that sum from observations.
Our finite periodic smoother is an illustration, not a new continuum approximation theorem. Its chosen boundary convention is part of the experiment.
\takeaway{Known-edge reconstruction removes the jump from the approximation problem. It does not remove the need to identify the jump.}

\vpage{16}{Why a Fourier cutoff cannot reproduce this jump pointwise}{A CONTINUOUS PULSE / EXACT ENDPOINT VALUES}{Chapter V: Section 6, Section 8.}
Use a separate period-one example on $[-1/2,1/2)$:
\[
j(t)=\mathbf 1_{\{|t|<1/4\}},\qquad
S_mj(t)=\tfrac12+\frac2\pi\sum_{n=1}^{m}\frac{\sin(\pi n/2)}{n}\cos(2\pi nt).
\]
This is the exact finite Fourier formula. The drawn curves sample it numerically; the endpoint calculation below is symbolic.
\vfig{fourier_jump}
\head{One value is fixed at every cutoff}
At $t=\pm1/4$, every summand vanishes because
$\sin(\pi n/2)\cos(\pi n/2)=0$ for integer $n$. Hence
\[
S_mj(\pm1/4)=\tfrac12 \qquad\text{for every }m\ge0.
\]
A larger cutoff resolves a narrower transition, but it does not make the finite trigonometric polynomial discontinuous. The two displayed cutoffs retain $11$ and $31$ signed frequency bins, respectively, not five and fifteen coefficients.
\head{An exact known-packet comparison}
Take $s(t)=\frac14\cos(2\pi t)+\frac1{10}\cos(6\pi t)$. For $m\ge3$ this smooth packet is already in the retained Fourier span. Thus $S_ms+j=s+j$ exactly in this special model, whereas $S_m(s+j)$ still replaces the pulse by $S_mj$.
The values assigned at the two endpoints matter for pointwise error, but do not affect an $L^p$ norm for finite $p$. The illustrated ringing is not a root enclosure or an asymptotic error certificate.
\takeaway{Increasing spectral resolution and retaining a known discontinuity are different operations, with different pointwise consequences.}

\vpage{17}{Small location errors, different norms}{MASK GEOMETRY / TRANSLATION BY ONE SIXTEENTH}{Chapter V: Theorem 7.1.}
In Lebesgue measure on the line, let
\[
\Omega=[\tfrac14,\tfrac34],\qquad
\widehat\Omega=[\tfrac5{16},\tfrac{13}{16}],\qquad a=\widehat a=3.
\]
The two intervals have the same length. Their disagreement consists of one missed strip and one extra strip, each of width $1/16$.
\vfig{mask_translation}
\[
|\Omega\mathbin{\triangle}\widehat\Omega|=\tfrac18,
\qquad j-\widehat j=3(\mathbf1_\Omega-\mathbf1_{\widehat\Omega}).
\]
\begin{center}
\begin{tabular}{c|ccc}
error measure & $L^1$ norm & squared $L^2$ norm & $L^\infty$ norm\\\hline
exact value & $3/8$ & $9/8$ & $3$
\end{tabular}
\end{center}
\head{Area shrinks while the missed jump keeps its height}
For a translation by $\delta$ with $0<|\delta|<1/2$, the symmetric difference has measure $2|\delta|$. At finite $p$,
\[
\|j-\widehat j\|_p=3(2|\delta|)^{1/p}.
\]
It tends to zero as the displacement tends to zero. The essential supremum stays $3$ for every nonzero displacement, because the error reaches height $3$ on a set of positive measure. At exactly zero displacement all norms vanish.
A nonempty disagreement set of measure zero would be different: it can affect the pointwise supremum while leaving the essential norm unchanged. This is why the paper distinguishes the two statements.
\takeaway{A shrinking boundary strip can give small integrated error and still contain a full-height jump error.}

\vpage{18}{Separate amplitude error from mask error}{COUNTING MEASURE / AN EXACT FINITE BUDGET}{Chapter V: Theorem 7.1.}
On the $16$ sample addresses, let $\Omega=\{6,\ldots,11\}$ and $\widehat\Omega=\{7,\ldots,12\}$. Use true amplitude $a=3$ and estimate $\widehat a=5/2$. There are six true samples, five overlapping samples, and two disagreement samples.
\vfig{error_budget}
The residual splits exactly as
\[
3\mathbf1_\Omega-\tfrac52\mathbf1_{\widehat\Omega}
=\tfrac12\mathbf1_\Omega+\tfrac52(\mathbf1_\Omega-\mathbf1_{\widehat\Omega}).
\]
The missed sample has error $3$, the five shared samples have error $1/2$, and the extra sample has error $-5/2$. With \emph{counting measure},
\[
\|j-\widehat j\|_1=8,
\qquad \|j-\widehat j\|_2^2=\tfrac{33}{2},
\qquad \|j-\widehat j\|_\infty=3.
\]
\head{Read the theorem as a budget, not an equality claim}
Here the $L^1$ bound is attained:
\[
|3-\tfrac52|\,|\Omega|+\tfrac52|\Omega\triangle\widehat\Omega|
=\tfrac12\cdot6+\tfrac52\cdot2=8.
\]
The $L^2$ triangle bound is $\frac12\sqrt6+\frac52\sqrt2$, and need not be sharp. A background estimate adds another term $\|s-\widehat s\|_p$; cancellation between terms can make the total smaller than the bound.
Normalized RMSE divides the $L^2$ norm by $\sqrt{16}$. A boundary-band RMSE uses the band size instead. Changing that denominator changes the reported quantity.
\takeaway{Mask error, amplitude error, and background error are separate mathematical inputs. A single reconstruction score conceals that separation.}

\vpage{19}{Geometry localizes error while data may leave it ambiguous}{A FINITE IMAGE / AN AMBIGUOUS OBSERVATION}{Chapter V: Section 7, Section 9.}
On the integer $8\times8$ grid, choose the disk-shaped mask
$\Omega=\{(i,j):(i-3)^2+(j-3)^2\le5\}$ and translate it one pixel to the right. Each mask contains $21$ pixels; their symmetric difference contains exactly $10$.
\vfig{mask_geometry}
\head{A geometric error calculation is not a segmentation algorithm}
For a common amplitude $a$, the squared counting-measure error is $10a^2$, entirely on the disagreement pixels. The full-image mean squared error is $10a^2/64$. Both are exact, but answer different normalization questions.
Now consider the four-sample observation
\[
y=(0,\tfrac32,3,3).
\]
It is consistent with either true signal
\[
u=(0,0,3,3),\qquad v=(0,3,3,3)
\]
under deterministic sup-norm noise at most $3/2$. Indeed $y=(u+v)/2$, and both error norms equal $3/2$. A rule receiving only $y$ and that noise bound cannot be guaranteed to select both distinct true masks correctly.
\head{What additional structure would have to do}
A lower noise level, separation, contrast, geometry, or a specified prior can change the identification problem. The example rules out a guarantee at the \emph{stated} noise radius and signal pair, not all recovery results. The manuscript's positive-contrast estimator is deliberately narrower than a universal segmentation method.
\takeaway{Exact synthesis can preserve an identified boundary. It cannot decide between two boundaries that the observation model makes indistinguishable.}

\vpage{20}{Two products, two packet laws}{POINTWISE SIGNALS VERSUS CAUCHY COEFFICIENTS}{Chapter V: Theorem 12.1.}
The word ``product'' must be typed. Multiplying two signal values at the same time address is not the same operation as multiplying generating series. The distinction changes which packets can interact.
\vfig{product_conventions}
\head{Pointwise multiplication preserves the same address}
For time-residue masks, $(P_{N,i}x)[n](P_{N,j}y)[n]=0$ when $i\ne j$. Therefore
\[
P_{N,k}(x\cdot y)=(P_{N,k}x)\cdot(P_{N,k}y).
\]
Here the dot means multiplication at the same sample, not convolution.
\head{Cauchy multiplication adds degrees}
For formal series $A(t),B(t)$, the coefficient of $t^n$ is a sum over $a+b=n$. Consequently the output residue is $i+j\pmod N$:
\[
P_{N,k}(AB)=\sum_{i+j\equiv k\ (N)}(P_{N,i}A)(P_{N,j}B).
\]
For $N=3$, take $A=1+2t+3t^2$ and $B=2+t^2$. Direct multiplication gives
\[
AB=2+4t+7t^2+2t^3+3t^4,
\qquad P_{3,0}(AB)=2+2t^3.
\]
The $2t^3$ term comes from residues $1+2$, not $0+0$. At power zero, the empty product contributes $1$ only in residue zero. These are formal coefficient identities and require no convergence assumption.
\takeaway{Before applying a packet multiplication rule, specify whether multiplication acts on sample values or on series coefficients.}

\vpage{21}{Nonlinear neighbors cross channel boundaries}{FULL INPUT TUPLES / EXACT CYCLIC INDEXING}{Chapter V: Section 12.}
For the periodic running signal define
\[
F(x)[n]=x[n]+\tfrac12x[n]x[n-1].
\]
This is nonlinear and not a convolution operator. Nevertheless its two-channel representation is exact, provided the neighbor comes from the correct channel and index.
\vfig{nonlinear_shift}
Put $e[q]=x[2q]$ and $o[q]=x[2q+1]$, with channel indices taken modulo eight. Then
\[
\begin{aligned}
F(x)[2q]&=e[q]+\tfrac12 e[q]o[q-1],\\
F(x)[2q+1]&=o[q]+\tfrac12 o[q]e[q].
\end{aligned}
\]
The even output needs the \emph{previous} odd sample; the odd output needs the current even sample. This is the same carry in the address system that appeared in the one-sample delay matrix.
\head{Check two actual outputs}
The running values are $x[1]=1/4$, $x[2]=1/2$, and $x[3]=1/4$, so
\[
F(x)[2]=\tfrac12+\tfrac12\tfrac12\tfrac14=\tfrac9{16},
\qquad F(x)[3]=\tfrac14+\tfrac12\tfrac14\tfrac12=\tfrac5{16}.
\]
Applying the same formula independently within each decimated array would use the wrong neighbors. Keeping the full channel tuple avoids that loss of information. The supplied suite compares every output sample in the two implementations using exact rational arithmetic.
No convergence, contraction, or approximation property follows merely from writing these exact coordinates. Those require hypotheses on the nonlinear map and the problem being solved.
\takeaway{Exact packet coordinates remain available for nonlinear maps. They do not turn cross-channel interactions into independent channel problems.}

\vpage{22}{One reflection circle, four very different impedances}{THE CAYLEY CHART / EXACT QUARTER TURNS}{Chapter V: Theorem 13.2.}
Use normalized impedance $z$ and reflection coefficient $\Gamma$:
\[
\Gamma=\frac{z-1}{z+1},\qquad z=\frac{1+\Gamma}{1-\Gamma}.
\]
The circle $|\Gamma|=1/2$ is a constant-mismatch orbit. Its four quarter-turn points map to the exact impedances shown below.
\vfig{cayley_orbit}
\begin{center}
\begin{tabular}{c|cccc}
$\Gamma$ & $1/2$ & $i/2$ & $-1/2$ & $-i/2$\\\hline
$z$ & $3$ & $3/5+4i/5$ & $1/3$ & $3/5-4i/5$
\end{tabular}
\end{center}
\head{The coordinates preserve structure, not Euclidean distances}
A quarter turn is a rotation in the reflection disk. After pulling back, it is a M\"obius map, not a Euclidean rotation of the impedance plane. Equal reflection magnitudes need not give equal impedance magnitudes. The opposite points $1/2$ and $-1/2$ become $3$ and $1/3$. The continuous impedance locus is the circle $(\operatorname{Re}z-5/3)^2+(\operatorname{Im}z)^2=16/9$, also drawn in the plate.
Finite passive impedances with $\operatorname{Re}z\ge0$ map to the closed disk except $\Gamma=1$. That missing point is obtained only by adjoining impedance infinity. The interior maps to the open right half-plane.
The classical Smith grid supplies the geometry; the cyclic labels specify exact rotations on it. Neither a new physical law nor a new optimal sampling theorem is implied by this coordinate description.
\takeaway{A simple action in one chart can be nonlinear in another. Conjugacy preserves the action law, not the apparent shape of a displacement.}

\vpage{23}{A pole is a chart boundary, not a broken group action}{HOMOGENEOUS COORDINATES / NO DIVISION BY ZERO}{Chapter V: Theorem 13.2, Corollary 13.3.}
Represent impedance by a nonzero pair $[u:v]$, where $[cu:cv]=[u:v]$ for $c\ne0$. A finite impedance is $[z:1]$; infinity is $[1:0]$. For a nonzero reflection multiplier $w$, use
\[
A(w)=\frac12\begin{pmatrix}1+w&1-w\\1-w&1+w\end{pmatrix}.
\]
\vfig{projective_poles}
\head{The action is defined before a ratio is taken}
The determinant is $w$, so the matrix cannot send a nonzero pair to zero. Direct multiplication gives
\[
A(w)A(w')=A(ww'),\qquad A(1)=I.
\]
For $w=i$, the orbit through zero is
\[
[0:1]\ \longmapsto\ [-i:1]\ \longmapsto\ [1:0]
\ \longmapsto\ [i:1]\ \longmapsto\ [0:1].
\]
The third point is a legitimate projective point, even though a finite impedance formula has a zero denominator there. The second iterate is $A(-1)$, which swaps the two coordinates. Thus a half turn exchanges $0$ and $\infty$ and sends each finite nonzero $z$ to $1/z$.
The matched point $\Gamma=0$, equivalently $z=1$, is fixed by every phase multiplier. It is not part of the boundary orbit above. The exact checks test homogeneous pairs and their rescaling equivalence, rather than introducing an artificial complex value for $1/0$.
\takeaway{Compute with homogeneous pairs first. A denominator becomes a chart-selection issue rather than an exception to the cyclic action.}

\vpage{24}{Residue labels are not impedance averages}{DYADIC--ODD FACTORIZATION / TWELVE ADDRESSES}{Chapter V: Proposition 14.1.}
At phase order $12=4\cdot3$, the map
\[
k\longmapsto(k\bmod4,\ k\bmod3)
\]
is a bijection from twelve residue addresses to the $4\times3$ grid. Every cell below contains its unique original address.
\vfig{crt_refinement}
\head{The factorization is an intersection of conditions}
For integer sample addresses, including negative ones,
\[
P_{12,k}=P_{4,k\bmod4}\,P_{3,k\bmod3}.
\]
The product on the right is composition of two residue masks: both congruences must hold. Coprimality makes their intersection a single class modulo twelve. It is not a formula for multiplying or adding impedances.
\head{Nested samples and nested projectors are different statements}
The four coarse reflection rays occur at indices $0,3,6,9$ of the twelve-ray grid. Thus the geometric sample set is nested. Separately, a coarse coefficient projector is the sum of its fine descendants in a common coefficient space.
There is no corresponding rule saying that a coarse impedance is the average of its descendants. The inverse Cayley map is nonlinear. Even opposite reflection points have impedances $3$ and $1/3$ in the preceding example; their average $5/3$ is not the impedance $1$ of their averaged reflection coefficient zero.
\takeaway{CRT factors addresses. Refinement conserves coefficient packets. Neither operation authorizes linear averaging through a nonlinear chart.}

\vpage{25}{Sensitivity grows before the pole is reached}{LOCAL DERIVATIVES / FINITE INCREMENTS}{Chapter V: Section 14.}
Away from $\Gamma=1$, the inverse Cayley map satisfies
\[
\frac{dz}{d\Gamma}=\frac2{(1-\Gamma)^2}.
\]
Along the real axis, three exact derivatives already show the scale change:
\[
\Gamma=\tfrac12:\ 8,\qquad
\Gamma=\tfrac9{10}:\ 200,\qquad
\Gamma=\tfrac{99}{100}:\ 20000.
\]
\vfig{sensitivity}
\head{A differential is not the full finite-error formula}
For a finite increment $\varepsilon$, with both denominators nonzero,
\[
z(\Gamma+\varepsilon)-z(\Gamma)
=\frac{2\varepsilon}{(1-\Gamma-\varepsilon)(1-\Gamma)}.
\]
At $\Gamma=9/10$ and $\varepsilon=1/100$, the actual increment is $20/9$, while derivative times increment is $2$. Both are calculated exactly. Their difference comes from using a differential in place of a finite increment. The derivative gives a first order approximation.
\head{What a proposed weight can and cannot certify}
Impedance is sensitive near $\Gamma=1$, and admittance near $\Gamma=-1$. The symmetric weight
\[
w(\Gamma)=|1-\Gamma|^{-2}+|1+\Gamma|^{-2}
\]
reflects both chart sensitivities. Its growth motivates denser sampling near the endpoints, but does not establish an optimal grid, interpolation rule, or measurement allocation. A finite-error estimate must also control how close the perturbed point comes to either pole.
\takeaway{Use the exact finite difference for an error budget. A large derivative diagnoses conditioning. Choosing an optimal algorithm requires further input.}

\vpage{26}{What the reader can now verify}{EVIDENCE GUIDE / THE FORMAL CONTRACT IS UNCHANGED}{Chapter V: Section 21 (formal correspondence), Section E15.}
\head{Three worked structures, three distinct kinds of conclusion}
\begin{tabularx}{\linewidth}{@{}p{.24\linewidth}X@{}}
\toprule
Structure & What is established in the original worked layer\\\midrule
Finite signal & Exact signed addressing, filter transfer, Haar reconstruction and energy, sparse coefficient error, and full-tuple nonlinear evaluation.\\
Finite packet tree & All $26$ complete prunings, actual projector sums, the unique cost-$9$ optimum, and a strict perturbation-margin certificate.\\
Reflection orbit & Exact rational-complex coordinates, a pole-safe projective action, CRT labels, and finite sensitivity calculations.\\\bottomrule
\end{tabularx}
\head{What each evidence layer means}
\textbf{Manuscript theorems.} The Part I formal correspondence maps twenty theorem-like statements to $225$ supporting Lean targets in $23$ proof modules. Its compiler and axiom records provide the retained execution evidence; the formal companion gives their scope and reproduction instructions. The additional Part II developments remain outside this twenty-result map.

\textbf{Recorded exact examples.} The symbolic visual suite was rerun for the recorded monograph review. All $181$ rational, algebraic, symbolic, and exhaustive finite checks passed. Seven deliberately false assertions were rejected. Numerical samples draw the continuous Fourier and sensitivity curves; they are not promoted to interval enclosures or universal error bounds.

\textbf{Recorded executions.} The fresh source-release checker passed 20 exact tests containing 305,933 assertions and passed 16 verifier tests, including positive and negative controls. The expanded-source numerical audit passed 11,052 finite checks, mixing exact and numerical calculations. The new extension examples separately passed 38 grouped exact checks and four false controls. These counts measure different operations and are not combined as theorem counts.\par
\textbf{Reproducibility.} The package preserves the original sources and their reports. Run \texttt{python3 build.py} here for an isolated PDF build. The extension finite checker is \texttt{checks/verify\_extension\_examples.py}; its contract and raw execution manifest are included. The original visual checker requires the listed optional dependencies. PDF assembly does not run Lean or either mathematical checker.
\head{The remaining boundaries are not decorative disclaimers}
Known-edge synthesis requires the correct packet. A finite cost optimum depends on the supplied objective. An infinite packet tree, a universal segmentation procedure, minimax recovery rates, and optimal Smith sampling require additional theorems. The core's broad edge-recovery problems are not resolved by these finite examples. Plate V.27 introduces the additional separated-jump recovery results and their explicit hypotheses.
The $16$-sample running signal is not the $1024$-sample manuscript experiment; the $26$-pruning example is not its $458{,}330$-pruning benchmark. Their smaller sizes allow every relevant number to be inspected exactly.
\takeaway{Coordinates, approximation, identification, optimization, and physical interpretation are separate layers. The examples make their interfaces visible without erasing their hypotheses.}

\vpage{27}{What the extended developments add}{PART II / REPRESENTATION BECOMES CONDITIONAL RECOVERY}{Chapter V: Section 21 (formal correspondence), Section E1, Section E15.}
The first twenty-six pages follow the corrected twenty-result formal core. The traditional paper's second part restores the broader mathematical developments of the earlier expanded manuscript. The next thirteen worked pages make their additional hypotheses and conclusions visible.
\head{Four transitions, with the missing input made explicit}
\begin{tabularx}{\linewidth}{@{}>{\raggedright\arraybackslash}p{.21\linewidth}YY@{}}\toprule
Transition & Additional input & Resulting control\\\midrule
Residues to lattices & A finite-index sublattice and one representative per coset & Exact addresses in several dimensions\\
Invertibility to stability & Positive lower frame bound and bounded perturbation & Quantified coefficient and reconstruction error\\
Known jumps to recovery & Separated jumps, sampling, a Lipschitz background, and a strict contrast gap & Exact jump cells and bounded location/amplitude error\\
Recovered packets to approximation & A bounded linear approximation operator and its norm & A complete two-term error estimate\\\bottomrule
\end{tabularx}
\head{The two parts have different formal evidence}
The inherited Lean map covers the twenty designated results of Part I through its 225 supporting targets. The additional traditional proofs in Part II are not silently included in that map. The new pages below are mathematical explanations and exact worked specializations of those developments, with a separate finite-example checker.
The original manuscript's broad edge-recovery questions remain broader than separated one-dimensional jumps. Conditions such as one jump per sample cell, an available wavelet support, and a positive detection margin do real work. In particular, the ambiguous observation on Plate V.19 still rules out an unrestricted recovery guarantee at its stated noise radius.
\head{How to read the extension}
Plates V.28--V.30 extend coordinates and stability. Plates V.31--V.34 identify what a singular packet contains and measure its error. Plates V.35--V.39 give three detection mechanisms. Plate V.40 closes the route by combining packet accuracy with a specified approximation operator.
\takeaway{A positive recovery theorem can coexist with an impossibility example. Their sampling models, noise bounds, and signal classes decide which conclusion applies.}

\vpage{28}{Six image phases are six exact address classes}{LATTICE COORDINATES / A RECTANGULAR EXAMPLE}{Chapter V: Theorem E3.2, Theorem E3.5, Remark E2.2.}
Take $L=2\ZZ\times3\ZZ$ inside $\ZZ^2$, with representatives
\[
R=\{0,1\}\times\{0,1,2\}.
\]
Every address has a unique expression $(i,j)=(r_1,r_2)+(2q_1,3q_2)$. The six channels are $x_r(2q_1,3q_2)=x(r_1+2q_1,r_2+3q_2)$.
\vfig[.80]{ext_lattice}
\platecaption{Each cell shows its residue label. Colors identify six disjoint classes; they do not average neighboring image pixels.}
\head{A negative address is still unique}
The point $(-1,7)$ has representative $(1,1)$ and lattice displacement $(-2,6)$, hence channel index $(-1,2)$. Floor division gives the correct answer in both coordinates.
For finite $p$, the channel norm uses the sum of the $p$th powers. At $p=\infty$, it uses their maximum. Both equal the norm on the original lattice because the cosets partition all addresses. In $\ell^\infty$, this coordinate identity does not imply that translates of an impulse are dense: their closure is only $c_0$.
\head{Refinement partitions a class again}
For $L'=4\ZZ\times6\ZZ\subset L$, the coarse representative $(1,2)$ has fine representatives
\[
(1,2),\ (3,2),\ (1,5),\ (3,5).
\]
Their original-lattice masks sum to the coarse mask. As on Plate V.4, decimated arrays are interleaved, not added entrywise.
\takeaway{Dimension changes the address tuple. Unique coset decomposition still supplies exact synthesis and the correct norm identity.}

\vpage{29}{A positive singular-value gap survives perturbation}{FRAME STABILITY / AN EXACT TWO-CHANNEL BANK}{Chapter V: Theorem E5.2, Theorem E5.7, Corollary E5.6.}
Consider the constant analysis bank
\[
E=\begin{pmatrix}1&0\\0&2\end{pmatrix},\qquad
E^*E=\operatorname{diag}(1,4).
\]
It has frame bounds $A=1$, $B=4$. For every perturbation with operator norm at most $\delta=1/4$, the squared-norm bounds remain at least $(1-1/4)^2=9/16$ and at most $(2+1/4)^2=81/16$.
\vfig[.87]{ext_frame}
\head{Both bounds can be attained}
Choose $\Delta E=\operatorname{diag}(-1/4,1/4)$. Then $\widetilde E=\operatorname{diag}(3/4,9/4)$ has precisely those two frame bounds. Its inverse has operator norm $4/3$. A coefficient error $\nu=(3/4,0)$ therefore produces signal error $(1,0)$, attaining the noise bound.
The strict condition $\delta<\sqrt A$ prevents the lower bound from vanishing. At $\delta=1$, the perturbation $\operatorname{diag}(-1,0)$ destroys the first channel entirely.
\head{An inverse and an approximate inverse answer different questions}
If the original synthesis $R=\operatorname{diag}(1,1/2)$ is retained after this perturbation, then
\[
R\widetilde E-I=\operatorname{diag}(-1/4,1/8).
\]
Its defect norm is $1/4$. Every input has relative reconstruction error at most that value, with equality on the first coordinate. For frequency-dependent banks the theorem uses an essential supremum on the unit circle, not a few sampled frequencies.
\takeaway{A lower singular-value gap quantifies stable recovery. A reconstruction defect quantifies the error of the actual synthesis bank being used.}

\vpage{30}{Nonunitary tree bounds propagate along branches}{LOCAL FRAME CONSTANTS / A MIXED-DEPTH CERTIFICATE}{Chapter V: Proposition E6.4.}
Suppose the root split has nonnegative squared-norm bounds $A=1/2$ and $B=2$. Retain the left child. Split the right child with bounds $3/4$ and $5/4$, then retain both grandchildren.
\vfig[.85]{ext_tree}
\head{Start with one at every retained node}
The left subtree has bounds $a_0=b_0=1$. The right subtree has $a_1=3/4$ and $b_1=5/4$. The recursive certificate gives
\[
a_{\rm root}=\tfrac12\min\{1,\tfrac34\}=\tfrac38,
\qquad
b_{\rm root}=2\max\{1,\tfrac54\}=\tfrac52.
\]
Consequently the total leaf coefficient energy lies between $3/8$ and $5/2$ times the root energy. These are guaranteed bounds; no assertion that every admissible collection of local maps attains both constants is needed.
\head{Why a minimum and maximum appear}
After the root split, its energy may distribute unevenly between the two children. The lower estimate must survive the least favorable child lower bound. The upper estimate must survive the largest child upper bound. Replacing these by an unweighted average would presume an energy distribution that the hypotheses do not supply.
When all local maps are unitary, every factor is one and the recursion recovers exact energy conservation. For general local maps, a positive lower bound alone establishes injectivity and stable inversion on the range; it does not by itself assert surjectivity onto every possible leaf tuple.
\takeaway{Exact tree coverage and good conditioning are separate inputs. The first chooses a coordinate route; the second bounds what that route does to norms.}

\vpage{31}{Derivative jumps have their own packets}{TRACE DATA / A STEP AND A RAMP}{Chapter V: Theorem E7.2, Corollary E7.3.}
On $[0,1]$, let $H_\tau=\mathbf1_{[\tau,1]}$ and take
\[
f(t)=t^2+2H_{1/3}(t)-3(t-2/3)_+.
\]
The function has a value jump $2$ at $1/3$ and a first-derivative jump $-3$ at $2/3$. Its packet and remainder are
\[
j_f(t)=2\phi_{0,1/3}(t)-3\phi_{1,2/3}(t),\qquad s_f(t)=t^2.
\]
\vfig[.92]{ext_trace}
\platecaption{The displayed curves are exact polynomial pieces. The solid jump is not connected across its discontinuity. The dashed curve is the global smooth remainder.}
\head{The normalization makes the trace coefficient literal}
The general packet is $\phi_{m,\tau}(t)=(t-\tau)_+^m/m!$. Its derivative of order $m$ has a unit trace jump, while its other classical trace jumps vanish. Thus its multiplier is exactly the jump being removed, with no hidden factorial.
Here the remainder has matching value and first-derivative traces and belongs to $W^{2,p}(0,1)$ for every $1\le p\le\infty$. Its second weak derivative is $2$. The singular packet accounts for the distributional boundary terms that would otherwise arise at the breakpoints.
\head{What the general theorem requires}
Each piece must lie in the stated Sobolev class and possess the specified one-sided traces. Subtracting the finite trace packet then joins those pieces into a global Sobolev remainder. This is a theorem about finitely many declared breakpoints, not a representation of every possible singular measure or arbitrary image boundary.
\takeaway{A value jump uses a step. A slope jump uses a ramp. The derivative order determines the packet and its normalization.}

\vpage{32}{The approximation operator also acts on packet error}{ESTIMATED PACKETS / AN EXACT AMPLIFICATION EXAMPLE}{Chapter V: Theorem E7.8.}
Let $Q$ be a bounded \emph{linear} operator and set $e=j-\widehat j$. Substituting $f=s+j$ into $\widehat f=Q(f-\widehat j)+\widehat j$ gives
\[
f-\widehat f=(I-Q)s+(I-Q)e.
\]
The second term includes the action of $Q$ on the packet error. Linearity is required for this split.
\vfig[.80]{ext_amplification}
\head{An oblique projection can increase the error}
In $\RR^2$, choose
\[
Q=\begin{pmatrix}1&2\\0&0\end{pmatrix},\qquad
s=(1,0),\qquad e=(0,1).
\]
Then $Q^2=Q$ and $Qs=s$, but $(I-Q)e=(-2,1)$. The smooth error is zero while the unit packet error becomes a reconstruction error of norm $\sqrt5$.
For a fully specified realization, take $j=e$ and $\widehat j=0$. Then $f=(1,1)$, $\widehat f=Qf=(3,0)$, and their difference is exactly $(-2,1)$.
\head{Use the actual operator bound}
The always-valid estimate in this setting is
\[
\|f-\widehat f\|\leq\|s-Qs\|+(1+\|Q\|)\|e\|.
\]
Here $\|Q\|=\sqrt5$, so the bound is conservative. For an orthogonal projection, $\|I-Q\|\le1$ gives a sharper estimate. The operator's geometry determines which bound can be used.
\takeaway{Recovering the packet and approximating the background are separate tasks, but the chosen approximation operator couples their final errors.}

\vpage{33}{An overlap matrix is an exact error ledger}{REGION VALUES / EVERY MISCLASSIFICATION HAS A WEIGHT}{Chapter V: Theorem E8.1, Theorem E8.4.}
Let two piecewise-constant signals have region values $a=(0,2,5)$ and $b=(0,2,5)$ on a counting-measure domain of twelve samples. Suppose their overlap counts are
\[
M=\begin{pmatrix}3&1&0\\0&2&1\\1&0&4\end{pmatrix}.
\]
The row totals $(4,3,5)$ are the true region sizes. The column totals are also $(4,3,5)$, although three samples are assigned to the wrong values.
\vfig[.80]{ext_overlap}
\head{Equal error counts need not mean equal reconstruction error}
The three off-diagonal entries contribute one sample each. Their value differences have magnitudes $2$, $3$, and $5$, so
\[
\|j-\widehat j\|_1=2+3+5=10,\qquad
\|j-\widehat j\|_2^2=4+9+25=38.
\]
The counting-measure supremum error is $5$. The full-domain mean squared error is $38/12=19/6$. A mismatch fraction of $3/12$ reports a different statistic and omits all three amplitude weights.
\head{The formula requires partitions}
For measurable partitions $(E_i)$ and $(F_j)$, with $M_{ij}=\mu(E_i\cap F_j)$,
\[
\|j-\widehat j\|_p^p=\sum_{i,j}|a_i-b_j|^pM_{ij},\qquad 1\le p<\infty.
\]
Each location belongs to exactly one intersection. Overlapping indicator packets that are summed at the same location must first be converted to their actual combined value regions before this partition formula is applied.
\takeaway{Geometry supplies the overlap counts. Signal values supply the weights. Their combination gives the norm relevant to reconstruction.}

\vpage{34}{A shifted step and a shifted ramp have different moduli}{LOCATION STABILITY / A ONE-SIXTEENTH DISPLACEMENT}{Chapter V: Corollary E8.2, Theorem E8.3, Theorem E8.7.}
Set $\tau=1/4$, $\widehat\tau=5/16$, and $\delta=1/16$. A unit step differs only on the interval between its two locations:
\[
\|H_\tau-H_{\widehat\tau}\|_2=\sqrt\delta=\tfrac14,
\qquad \|H_\tau-H_{\widehat\tau}\|_\infty=1.
\]
\vfig[.89]{ext_location}
\head{For a ramp, the error height shrinks too}
Write $r_\tau(t)=(t-\tau)_+$. The difference grows linearly from zero to $\delta$ on $[\tau,\widehat\tau]$ and stays at $\delta$ on $[\widehat\tau,1]$. Direct integration gives
\[
\|r_\tau-r_{\widehat\tau}\|_2^2
=\frac{\delta^3}{3}+(1-\widehat\tau)\delta^2
=\frac{17}{6144}<\frac1{256}=\delta^2.
\]
Thus the $L^2$ ramp error is at most $\delta$, and its supremum error is exactly $\delta$. Positive-order truncated powers are Lipschitz in location, with the order-dependent constants supplied by the theorem.
\head{The geometric analogue needs a tube hypothesis}
If $E\triangle F\subseteq N_\delta(\partial E)$ and the tube measure is bounded by $C\delta$ for the stated range $0<\delta<\delta_0$, then an equal-amplitude packet has finite-$p$ error at most $|a|(C\delta)^{1/p}$. A mere visual impression that two boundaries are close does not establish the inclusion or the tube estimate.
\takeaway{A displaced value jump keeps its full height on a smaller set. A displaced derivative-jump packet also decreases in height.}

\vpage{35}{A strict contrast gap locates every separated jump cell}{DETERMINISTIC RECOVERY / A NUMERICAL CERTIFICATE}{Chapter V: Lemma E9.1, Theorem E9.2, Example E9.7.}
Assume a real signal is a Lipschitz background plus jumps. No jump occurs at a sample point, and each grid cell contains at most one jump. With spacing $h$, background constant $M$, and sample noise bounded by $\varepsilon$, write
\[
b_h=Mh+2\varepsilon.
\]
A no-jump contrast has magnitude at most $b_h$. A jump contrast differs from its true amplitude by at most $b_h$.
\vfig[.89]{ext_threshold}
\head{Use the manuscript's exact numerical inputs}
For $h=1/256$, $M=1/2$, $\varepsilon=1/500$, and $a_{\min}=2/25$,
\[
b_h=\frac{381}{64000}<\frac3{500},\qquad
2b_h<\frac2{25}.
\]
The threshold $\theta=1/25$ lies strictly between $b_h$ and $a_{\min}-b_h$. Thresholding $|y_{k+1}-y_k|$ therefore identifies exactly the jump-containing cells. Their midpoints locate each jump within $1/512$, and each estimated amplitude differs by at most $381/64000$.
\head{The recovered object is a cell, then an estimate}
The theorem does not claim the exact real-valued breakpoint from a finite grid. It supplies the correct cell and a controlled representative of that cell. For real jumps, the amplitude error is too small to reverse the sign. Two opposite jumps in one cell could cancel in the contrast, which is why the separation assumption cannot be dropped.
For $Q$ jumps the finite-$p$ packet bound contains both $Qb_h$ and the location term $(\tfrac h2\sum_q|a_q|^p)^{1/p}$. A growing number of jumps or shrinking separation requires its own analysis.
\takeaway{Exact detection means exact cell identification under a strict margin. Location and amplitude remain quantitative estimates with explicit errors.}

\vpage{36}{Repeated observations shrink the noise side of the gap}{SUBGAUSSIAN RECOVERY / A CONSERVATIVE SAMPLE BUDGET}{Chapter V: Theorem E9.8, Corollary E9.9, Corollary E9.11.}
Suppose all noise observations are independent and centered, with moment-generating bound $\mathbb E e^{\lambda\eta}\le e^{\sigma^2\lambda^2/2}$. Averaging $R$ observations at each of $n+1$ sample sites replaces $\sigma$ by $\sigma/\sqrt R$. A sufficient recovery condition is
\[
a_{\min}>2Mh+\frac{4\sigma}{\sqrt R}\sqrt{\log(2n/\delta)}.
\]
The probability of simultaneous recovery is then at least $1-\delta$, subject to the same separated-jump sampling model.
\vfig[.86]{ext_repetition}
\head{One explicit conservative choice}
Take $n=256$, $M=1/2$, $\sigma=1/50$, $\delta=1/20$, and $a_{\min}=2/25$. Since $\log(10240)<10<(16/5)^2$, choosing $R=16$ makes the right-hand side strictly less than
\[
\frac1{256}+\frac1{50}\frac{16}{5}
=\frac{2173}{32000}<\frac2{25}.
\]
Thus sixteen independent observations per sample site suffice for this bound. It is a sufficient budget, not a claim of optimal sample complexity. The exact checker certifies $e^{10}>10240$ using a rational Taylor partial sum, so the displayed logarithm bound is not based on rounded numerical logarithms.
\head{Dependence between adjacent contrasts is allowed}
Neighboring differences share one sample. The proof does not require those differences to be independent: it bounds each tail and applies a union bound across all cells. Independence of the underlying sample noises supplies the stated sub-Gaussian difference proxy.
For a fixed separated signal model, a noise scale satisfying $\sigma_n\sqrt{\log n}\to0$ gives the stated recovery consistency in probability. A fixed positive noise scale with one observation per site does not satisfy that condition.
\takeaway{More grid points and more repeated observations affect different terms. Finer spacing reduces background variation; averaging reduces noise.}

\vpage{37}{Haar sees a jump only through its position in the support}{SCALE SEPARATION / A QUANTIFIED THRESHOLD}{Chapter V: Theorem E10.1, Theorem E10.2, Corollary E10.3.}
Use a normalized Haar wavelet on an available interval of length $L$ on the real line, positive on its left half and negative on its right half. Here $H_\tau$ is the half-line step $\mathbf1_{[\tau,\infty)}$, so there is no artificial second endpoint jump.
If $u$ is the relative jump position in the support, then
\[
|\langle aH_\tau,\psi\rangle|
=|a|\sqrt L\min\{u,1-u\}.
\]
The coefficient vanishes when the jump lies exactly at a support boundary.
\vfig[.86]{ext_haar_margin}
\head{A worked gap with all three contributions}
Take $L=1/4$, $M=1$, noise bound $\varepsilon=1/100$, jump size $|a|=1$, and relative boundary distance at least $\rho=1/4$. The smooth and noise contributions together are bounded by
\[
\frac M4L^{3/2}+\varepsilon\sqrt L
=\frac1{32}+\frac1{200}=\frac{29}{800}.
\]
The jump contributes at least $|a|\rho\sqrt L=1/8$. The observed coefficient therefore has magnitude at least $71/800$. The threshold $1/16=50/800$ separates every such jump support from every support containing no jump.
\head{The available support is part of the hypothesis}
A fixed extension is needed if the observation domain ends inside a desired wavelet support. Also, the selected support must contain only the specified jump. An unmodeled second jump can reinforce or cancel the coefficient.
\takeaway{Scale determines the smooth contribution, but alignment determines the jump contribution. A detection guarantee needs both.}

\vpage{38}{Two half-shifted grids remove exact alignment failure}{HAAR SUPPORT GEOMETRY / A QUARTER-CELL GUARANTEE}{Chapter V: Theorem E10.4, Corollary E10.5.}
Normalize the support length to $L=1$. The first grid has boundaries at integers; the second has boundaries at half-integers. For a point at position $u$ modulo one, consider its distance to the nearest boundary of each grid.
\vfig[.89]{ext_two_grids}
\head{At least one distance is one quarter or larger}
If $1/4\le u\le3/4$, the original cell works. Outside that range, a shifted cell works. At $u=1/8$, the original margin is $1/8$ but the shifted cell $[-1/2,1/2)$ gives margin $3/8$. At $u=1/4$, both margins equal $1/4$, showing that the uniform constant cannot be increased for these two grids.
Restoring $L$, every jump has a support with boundary distance at least $L/4$ in one grid. If $L$ is smaller than the minimum jump separation, that support contains no other jump.
\head{Insert the geometric guarantee into the analytic estimate}
The alignment-independent gap condition is
\[
\frac{|a|}{4}>2\varepsilon+\frac M2L.
\]
For the numbers on Plate V.37, the right side is $1/50+1/8=29/200$, strictly below $1/4$. Hence one of the two grids supplies a detectable support for every jump, regardless of its position relative to the original grid.
This certifies a useful support per jump. It does not by itself specify a unique rule for merging duplicate detections from two grids into exact breakpoint estimates.
\takeaway{A second grid repairs a geometric blind spot. The contrast inequality and separation hypotheses still supply the recovery guarantee.}

\vpage{39}{Higher differences annihilate polynomial interiors}{BOUNDARY STENCILS / SIGNAL RESPONSE AND NOISE GAIN}{Chapter V: Theorem E10.6, Proposition E10.7, Corollary E10.8.}
Let $x(n)=n+2(n-3)_+$ on the integer line. Its slope changes from one to three at the breakpoint $3$. The second forward difference is
\[
\Delta^2x(n)=x(n+2)-2x(n+1)+x(n).
\]
\vfig[.87]{ext_differences}
\head{An exact small stencil calculation}
The values at $n=0,\ldots,6$ are $(0,1,2,3,6,9,12)$. The second differences at $n=0,\ldots,4$ are $(0,0,2,0,0)$. The only nonzero stencil straddles the slope change.
For a piecewise polynomial of degree at most $d$, $\Delta^{d+1}$ vanishes whenever its complete stencil lies inside one polynomial piece. Candidate responses are confined to at most $d+1$ starting addresses before each breakpoint. They are boundary bands, not necessarily exact breakpoint locations.
\head{Differencing also acts on noise}
The absolute binomial coefficients of $\Delta^m$ sum to $2^m$. With $\|\eta\|_\infty\le1/16$, second-difference noise is bounded by $4/16=1/4$. The true response $2$ remains at least $7/4$, while all responses away from the breakpoint band remain at most $1/4$. Threshold one separates them.
The general robust-band statement assumes that each breakpoint has at least one sufficiently large true response. A breakpoint can fail that assumption through its amplitude or local configuration. Annihilation away from boundaries does not prove a lower bound at every boundary.
\takeaway{A vanishing interior response is structural. Reliable detection also requires the boundary response to exceed the amplified noise.}

\vpage{40}{Conditional Gibbs control closes the reconstruction estimate}{THE TWO ERRORS / A COMPLETE READING CONTRACT}{Chapter V: Theorem E11.2, Theorem E11.3.}
The exact retained-packet approximation shares the same discontinuity as its target. If $f=s+j$ and $Q_ms\to s$ uniformly, then
\[
\|f-(Q_ms+j)\|_\infty=\|s-Q_ms\|_\infty\longrightarrow0.
\]
The approximants themselves may be discontinuous. There is no conflict with the failure of continuous Fourier polynomials to converge uniformly to a jump.
\begin{center}\includegraphics[width=.78\linewidth]{visual_graphics/annals_gibbs_comparison.pdf}\end{center}
\platecaption{Preserved vector figure from the expanded source. It illustrates the choice of approximation target; it is not an independent universal recovery certificate.}
\head{Estimated packets need the operator norm too}
If the bounded linear operators satisfy $\sup_m\|Q_m\|_{L^p\to L^p}\le C_Q$, then
\[
\|f-\widehat f_m\|_p
\le\|s-Q_ms\|_p+(1+C_Q)\|j-\widehat j_m\|_p.
\]
For an orthogonal $L^2$ projection one may use $\|I-Q_m\|\le1$, sharpening the packet multiplier to one. Such a bound must be justified in the norm actually being used. In particular, uniform boundedness of Fourier partial sums in the supremum norm is not supplied by this theorem.
\head{What the complete guide establishes}
Plates V.1--V.26 preserve the original worked layer. Plates V.27--V.40 explain the additional traditional developments, outside the inherited twenty-result Lean map. The new exact examples have their own finite arithmetic record. They do not extend the formal coverage merely by appearing beside formally mapped results.
The final package keeps recorded computations and inherited compiler reports attached to their stated scope. The symbolic visual suite was rerun, with 181 checks passed and seven invalid controls rejected. Lean compiler evidence remains inherited. No fresh Lean run is claimed.
\takeaway{A complete error claim specifies the signal class, the recovered packet, the approximation operator, and the norm. Each term can then be improved and checked on its own terms.}
